%Daten zur Institution

%\input{Dozentdaten}

%\renewcommand{\fachbereich}{Fachbereich}

%\renewcommand{\dozent}{Prof. Dr. . }

%Klausurdaten

\renewcommand{\klausurgebiet}{ }

\renewcommand{\klausurtyp}{ }

\renewcommand{\klausurdatum}{ . 20}

\klausurvorspann {\fachbereich} {\klausurdatum} {\dozent} {\klausurgebiet} {\klausurtyp}

%Daten für folgende Punktetabelle

\renewcommand{\aeins}{  3  }

\renewcommand{\azwei}{  3   }

\renewcommand{\adrei}{  3  }

\renewcommand{\avier}{  0  }

\renewcommand{\afuenf}{  2  }

\renewcommand{\asechs}{  0  }

\renewcommand{\asieben}{  3  }

\renewcommand{\aacht}{  7  }

\renewcommand{\aneun}{  0  }

\renewcommand{\azehn}{  7  }

\renewcommand{\aelf}{  4  }

\renewcommand{\azwoelf}{  3  }

\renewcommand{\adreizehn}{  3   }

\renewcommand{\avierzehn}{  7  }

\renewcommand{\afuenfzehn}{  2  }

\renewcommand{\asechzehn}{  9  }

\renewcommand{\asiebzehn}{  7  }

\renewcommand{\aachtzehn}{  63  }

\renewcommand{\aneunzehn}{    }

\renewcommand{\azwanzig}{    }

\renewcommand{\aeinundzwanzig}{    }

\renewcommand{\azweiundzwanzig}{    }

\renewcommand{\adreiundzwanzig}{    }

\renewcommand{\avierundzwanzig}{    }

\renewcommand{\afuenfundzwanzig}{    }

\renewcommand{\asechsundzwanzig}{    }

\punktetabellesiebzehn

\klausurnote

\newpage

\setcounter{section}{0}

__NOEDITSECTION__

\inputaufgabeklausurloesung
{3}

{

Definiere die folgenden
\zusatzklammer {kursiv gedruckten} {} {} Begriffe.
\aufzaehlungsechs{Die
\stichwort {Gauß-Abbildung} {}
zu einer differenzierbaren Hyperfläche

\mavergleichskette
{\vergleichskette
{ Y
}
{ \subseteq }{ \R^n
}
{  }{
}
{    }{
}
{   }{
}
}
{}{}{.}

}{Eine
\stichwort {topologische Karte} {}
auf einer
\definitionsverweis {topologischen Mannigfaltigkeit}{}{}
$M$.

}{Eine  $C^1$-\stichwort {differenzierbare Mannigfaltigkeit} {} $M$.

}{Die $n$-te \stichwort {äußere Potenz} {} zu einem
$K$-\definitionsverweis {Vektorraum}{}{}
$V$
\zusatzklammer {es genügt, das Symbol dafür anzugeben} {} {.}

}{Eine
\stichwort {lokale Isometrie} {}
\maabb {\varphi} {L} {M
} {}
zwischen riemannschen Mannigfaltigkeiten.

}{Ein
\stichwort {lokal integrabler} {}
Zusammenhang auf einem differenzierbaren Vektorbündel
\maabb {p} {E} {X
} {}
über einer differenzierbaren Mannigfaltigkeit $X$.
}

}
{

\aufzaehlungsechs{Es sei eine
\definitionsverweis {Orientierung}{}{}
auf $Y$ fixiert. Dann heißt die Abbildung
\maabbdisp {} {Y} { S^{n-1}
} {,}
die jeden Punkt

\mavergleichskette
{\vergleichskette
{ P
}
{ \in }{ Y
}
{  }{
}
{    }{
}
{   }{
}
}
{}{}{}
auf seinen durch die Orientierung fixierten Einheitsnormalenvektor abbildet, die
Gauß-Abbildung
zu $Y$.
}{Eine topologische Karte ist jede
\definitionsverweis {Homöomorphie}{}{}
\maabbdisp {\varphi} {U} {V
} {,}
wobei

\mavergleichskette
{\vergleichskette
{ U
}
{ \subseteq }{ M
}
{  }{
}
{    }{
}
{   }{
}
}
{}{}{}
und

\mavergleichskette
{\vergleichskette
{ V
}
{ \subseteq }{ \R^n
}
{  }{
}
{    }{
}
{   }{
}
}
{}{}{}
\definitionsverweis {offen}{}{}
sind.
}{Ein
\definitionsverweis {topologischer}{}{}
\definitionsverweis {Hausdorffraum}{}{}
$M$ zusammen mit einer
\definitionsverweis {offenen Überdeckung}{}{}

\mathl{M= \bigcup_{i \in I} U_i}{} und
\definitionsverweis {Karten}{}{}
\maabbdisp {\alpha_i} {U_i } { V_i
} {}
mit
\mathl{V_i \subseteq \R^n}{} offen derart, dass die
\definitionsverweis {Übergangsabbildungen}{}{}
\maabbdisp {\alpha_j \circ (\alpha_i)^{-1}} {V_i \cap \alpha_i(U_i \cap U_j)   } { V_j \cap \alpha_j(U_i \cap U_j)
} {}
$C^1$-\definitionsverweis {Diffeomorphismen}{}{}
für alle
\mathl{i,j \in  I}{} sind, heißt differenzierbare Mannigfaltigkeit.
}{Man nennt den $K$-Vektorraum
\mathl{\bigwedge^n V}{} das $n$-te Dachprodukt von $V$.
}{Eine
\definitionsverweis {differenzierbare Abbildung}{}{}
\maabb {\varphi} {L} {M
} {}
heißt
lokale Isometrie,
wenn für jeden Punkt

\mavergleichskette
{\vergleichskette
{ P
}
{ \in }{ L
}
{  }{
}
{    }{
}
{   }{
}
}
{}{}{}
die
\definitionsverweis {Tangentialabbildung}{}{}
\maabbdisp {T_P \varphi} {T_PL} { T_{\varphi(P) } M
} {}
eine
\definitionsverweis {Isometrie}{}{}
bezüglich der gegebenen Skalarprodukte ist.
}{Der Zusammenhang heißt
lokal integrabel,
wenn es zu jedem Punkt

\mavergleichskette
{\vergleichskette
{ e
}
{ \in }{ E
}
{  }{
}
{    }{
}
{   }{
}
}
{}{}{}
einen auf einer offenen Umgebung

\mavergleichskette
{\vergleichskette
{ p(e)
}
{ \in }{ U
}
{ \subseteq }{ X
}
{    }{
}
{   }{
}
}
{}{}{}
definierten
\definitionsverweis {horizontalen Schnitt}{}{}
\maabbdisp {s} {U} {E \vert_U
} {}
durch $e$ gibt.
}

 }

__NOEDITSECTION__

\inputaufgabeklausurloesung
{3}

{

Formuliere die folgenden Sätze.
\aufzaehlungdrei{Der
\stichwort {Isometriesatz für den Paralleltransport auf einer Hyperfläche} {}

\mavergleichskette
{\vergleichskette
{  Y
}
{ \subseteq }{  \R^n
}
{  }{
}
{    }{
}
{   }{
}
}
{}{}{.}}{Das
\stichwort {Theorema egregium} {.}}{Der Satz über die Partition der Eins.}

}
{

\aufzaehlungdrei{\faktsituation {Es sei

\mavergleichskette
{\vergleichskette
{ Y
}
{ \subseteq }{ W
}
{ \subseteq }{ \R^n
}
{    }{
}
{   }{
}
}
{}{}{,}
$W$
\definitionsverweis {offen}{}{,}
eine
\definitionsverweis {differenzierbare Hyperfläche}{}{}
und sei
\maabb {\gamma} {[a,b]} {Y
} {}
eine
\definitionsverweis {differenzierbare Kurve}{}{}
mit

\mavergleichskette
{\vergleichskette
{ \gamma(a)
}
{ = }{ P
}
{  }{
}
{    }{
}
{   }{
}
}
{}{}{}
und

\mavergleichskette
{\vergleichskette
{ \gamma(b)
}
{ = }{ Q
}
{  }{
}
{    }{
}
{   }{
}
}
{}{}{.}}
\faktfolgerung {Dann ist der
\definitionsverweis {Paralleltransport}{}{}
längs $\gamma$ eine
\definitionsverweis {Isometrie}{}{}
\maabbdisp {} {T_PY} {T_QY
} {.}}
\faktzusatz {}
\faktzusatz {}}{\faktsituation {Es sei

\mavergleichskette
{\vergleichskette
{ Y
}
{ \subseteq }{ \R^3
}
{  }{
}
{    }{
}
{   }{
}
}
{}{}{}
eine
\definitionsverweis {orientierte Fläche}{}{}
und sei
\maabbdisp {\varphi} { V } { Y
} {,}

\mavergleichskette
{\vergleichskette
{ V
}
{ \subseteq }{ \R^2
}
{  }{
}
{    }{
}
{   }{
}
}
{}{}{,}
eine zweifach differenzierbare lokale Parametrisierung von $Y$ mit den Parametern $u,v$. Es sei
\mathl{\left(  g_{ij}  \right)}{} die
\definitionsverweis {erste Fundamentalmatrix}{}{}
auf $V$.}
\faktfolgerung {Dann gilt für die
\definitionsverweis {Gaußsche Krümmung}{}{}
unter Verwendung der
\definitionsverweis {Christoffelsymbole}{}{}
die Beziehung

\mavergleichskettedisphandlinks
{\vergleichskettedisphandlinks
{ K
}
{ =} { { \frac{   1  }{ g_{11}  } }  \left(  \partial_2 \Gamma_{11}^2 - \partial_1 \Gamma_{12}^2 + \Gamma^1_{11} \Gamma^2_{12}  + \Gamma_{11}^2 \Gamma_{22}^2 - \Gamma_{12}^1 \Gamma_{11}^2-  \left(  \Gamma_{12}^2  \right)^2  \right)
}
{ } {
}
{ } {
}
{ } {
}
}
{}{}{.}}
\faktzusatz {}
\faktzusatz {}}{Es sei $M$ eine
\definitionsverweis {differenzierbare Mannigfaltigkeit}{}{}
mit einer
\definitionsverweis {abzählbaren Basis}{}{}
der Topologie. Dann gibt es zu jeder offenen Überdeckung eine der Überdeckung untergeordnete
\definitionsverweis {stetig differenzierbare}{}{}
\definitionsverweis {Partition der Eins}{}{.}}

 }

__NOEDITSECTION__

\inputaufgabeklausurloesung
{3}

{

Es seien
\maabbdisp {f,g} {\R} { \R^n
} {}
\definitionsverweis {differenzierbare Kurven}{}{.}
Berechne die
\definitionsverweis {Ableitung}{}{}
der Funktion
\maabbeledisp {} {\R} {\R
} { t } { \left\langle f(t) , g(t)  \right\rangle
} {.}
Formuliere das Ergebnis mit dem Skalarprodukt.

}
{

Es seien $f_i(t)$ bzw. $g_i(t)$ die Komponentenfunktionen von
\mathkor {} {f} {bzw.} {g} {.}
Die in Frage stehende Funktion ist

\mavergleichskettedisp
{\vergleichskette
{  \left\langle f(t) , g(t)  \right\rangle
}
{ =} { \sum_{i = 1}^n  f_i(t) g_i(t)
}
{ } {
}
{ } {
}
{ } {
}
}
{}{}{}
mit der Ableitung
\zusatzklammer {auf jeden Summanden wendet man die Produktregel an} {} {}

\mathdisp {\sum_{i = 1}^n  f'_i(t) g_i(t)  + \sum_{i = 1}^n  f_i(t) g'_i(t)} { . }

Mit dem Skalarprodukt kann man dies als

\mathdisp {\left\langle f'(t) ,  g(t)  \right\rangle + \left\langle f(t) ,  g'(t)  \right\rangle} {  }

schreiben.

 }

__NOEDITSECTION__

\inputaufgabeklausurloesung
{0}

{

}
{  }

__NOEDITSECTION__

\inputaufgabeklausurloesung
{2}

{

Es sei

\mavergleichskette
{\vergleichskette
{ Y
}
{ \subseteq }{ W
}
{ \subseteq }{ \R^n
}
{    }{
}
{   }{
}
}
{}{}{,}
$W$
\definitionsverweis {offen}{}{,}
eine
\definitionsverweis {differenzierbare Hyperfläche}{}{}
und sei
\maabb {\gamma} {I =[a,b]} {Y
} {}
eine
\definitionsverweis {geodätische Kurve}{}{.}
Zeige, dass der
\definitionsverweis {Paralleltransport}{}{}
längs $\gamma$ den tangentialen Vektor

\mavergleichskette
{\vergleichskette
{ \gamma'(a)
}
{ \in }{ T_{\gamma(a)}Y
}
{  }{
}
{    }{
}
{   }{
}
}
{}{}{}
in den tangentialen Vektor

\mavergleichskette
{\vergleichskette
{ \gamma'(b)
}
{ \in }{ T_{\gamma(b)}Y
}
{  }{
}
{    }{
}
{   }{
}
}
{}{}{}
überführt.

}
{

Nach
[[Differenzierbare Hyperfläche/Differenzierbare Kurve/Ableitung parallel/Geodätische/Fakt|Lemma 6.8 (Differentialgeometrie (Osnabrück 2023))]]
ist das Geschwindigkeitsfeld $\gamma'$ ein
\definitionsverweis {paralleles Vektorfeld}{}{}
längs $\gamma$. Daher überführt der
\definitionsverweis {Paralleltransport}{}{}
$\Psi_\gamma$  zu $\gamma$ den Vektor

\mavergleichskette
{\vergleichskette
{ \gamma'(a)
}
{ \in }{  T_{\gamma(a)} Y
}
{  }{
}
{    }{
}
{   }{
}
}
{}{}{}
in den Vektor

\mavergleichskette
{\vergleichskette
{ \gamma'(b)
}
{ \in }{  T_{\gamma(b)} Y
}
{  }{
}
{    }{
}
{   }{
}
}
{}{}{}
über.

 }

__NOEDITSECTION__

\inputaufgabeklausurloesung
{0}

{

}
{  }

__NOEDITSECTION__

\inputaufgabeklausurloesung
{3}

{

Ist die Abbildung
\maabbeledisp {\varphi} {S^1} {S^1
} {(x,y)} {( \betrag {  x  } ,y)
} {,}
\definitionsverweis {differenzierbar}{}{?}

}
{

Wir betrachten die Verknüpfung von Abbildungen

\mathdisp {\R \stackrel{ \theta }{\longrightarrow} S^1 \stackrel{\varphi}{\longrightarrow} S^1 \subset \R^2 \stackrel{p_1}{\longrightarrow} \R} { , }

wobei

\mavergleichskettedisp
{\vergleichskette
{ \theta(t)
}
{ =} {(\cos  t , \sin  t )
}
{ } {
}
{ } {
}
{ } {
}
}
{}{}{}
und $p_1$ die erste Projektion ist. Die trigonometrische Parametrisierung, die Inklusion
\zusatzklammer {einer abgeschlossenen Untermannigfaltigkeit} {} {}
und die Projektion sind differenzierbar. Wäre $\varphi$ differenzierbar, so müsste auch die Gesamtabbildung differenzierbar sein. Diese ist aber
\maabbeledisp {} {\R} {\R
} { t } {\betrag {  \cos  t  }
} {.}
Diese ist an der Stelle
\mathl{t= \pi/2}{} nicht differenzierbar, da dort die linksseitige Steigung $-1$ und die rechtsseitige Steigung $1$ ist. Die Abbildung $\varphi$ ist also nicht differenzierbar.

 }

__NOEDITSECTION__

\inputaufgabeklausurloesung
{7}

{

Es sei $M$ eine kompakte differenzierbare Mannigfaltigkeit mit einer stetigen positiven Volumenform $\omega$. Zeige, dass

\mavergleichskettedisp
{\vergleichskette
{ \int_M \omega
}
{ <} { \infty
}
{ } {
}
{ } {
}
{ } {
}
}
{}{}{}
ist.

}
{

Zu jedem Punkt

\mavergleichskette
{\vergleichskette
{ P
}
{ \in }{ M
}
{  }{
}
{    }{
}
{   }{
}
}
{}{}{}
gibt es eine offene Kartenumgebung

\mavergleichskette
{\vergleichskette
{ P
}
{ \in }{ U
}
{  }{
}
{    }{
}
{   }{
}
}
{}{}{}
und eine Kartenabbildung
\maabbdisp {\alpha} { U } { V
} {}
mit

\mavergleichskette
{\vergleichskette
{ V
}
{ \subseteq }{ \R^n
}
{  }{
}
{    }{
}
{   }{
}
}
{}{}{}
offen und so, dass

\mavergleichskette
{\vergleichskette
{ \alpha^{-1 *} \omega
}
{ = }{ fdx_1 \wedge \ldots \wedge dx_n
}
{  }{
}
{    }{
}
{   }{
}
}
{}{}{}
ist mit
\mathl{f}{} stetig und positiv. Wir finden auch eine offene Umgebung

\mavergleichskette
{\vergleichskette
{ P
}
{ \in }{ U'
}
{ \subseteq }{ U
}
{    }{
}
{   }{
}
}
{}{}{,}
die homöomorph zu einem offenen Ball

\mavergleichskette
{\vergleichskette
{ U'
}
{ \cong }{ B'
}
{ \subseteq }{ V
}
{    }{
}
{   }{
}
}
{}{}{}
ist, wobei man auch annehmen kann, dass der Abschluss des Balles ganz in $V$ liegt. Der abgeschlossene Ball ist abgeschlossen und beschränkt, daher ist die stetige Funktion $f$ darauf und somit auch auf $U'$ beschränkt. Es folgt, dass
\mathl{\int_{U'} \omega}{} endlich ist, wobei $U'$ eine offene Umgebung von $P$ ist.

Diese offenen Mengen

\mavergleichskette
{\vergleichskette
{ U'
}
{ = }{ U'(P)
}
{  }{
}
{    }{
}
{   }{
}
}
{}{}{}
überdecken $M$. Wegen der Kompaktheit gibt es eine endliche Überdeckung

\mavergleichskettedisp
{\vergleichskette
{ M
}
{ =} { \bigcup_{i = 1} ^n U_i
}
{ } {
}
{ } {
}
{ } {
}
}
{}{}{}
mit

\mavergleichskettedisp
{\vergleichskette
{ \int_{U_i } \omega
}
{ <} { \infty
}
{ } {
}
{ } {
}
{ } {
}
}
{}{}{.}
Wegen der Positivität gilt somit

\mavergleichskettedisp
{\vergleichskette
{ \int_{M} \omega
}
{ \leq} { \sum_{i = 1}^n  \left(  \int_{U_i } \omega  \right)
}
{ <} { \infty
}
{ } {
}
{ } {
}
}
{}{}{.}

 }

__NOEDITSECTION__

\inputaufgabeklausurloesung
{0}

{

}
{  }

__NOEDITSECTION__

\inputaufgabeklausurloesung
{7 (1+2+4)}

{

\aufzaehlungdreiabc{ Man gebe eine topologische Eigenschaft an, die zeigt, dass das
\definitionsverweis {offene Einheitsintervall}{}{}
und die
\definitionsverweis {Kreissphäre}{}{}
$S^1$ nicht
\definitionsverweis {homöomorph}{}{}
sind.
}{Zeige, dass jede stetige
$1$-\definitionsverweis {Differentialform}{}{}
auf
\mathl{]0,1[}{}
\definitionsverweis {exakt}{}{}
ist.
}{ Man gebe eine konkrete
\definitionsverweis {geschlossene}{}{}
$1$-Form auf $S^1$ an, die nicht exakt ist.
}

}
{

a) Nach
[[Kompaktheit/Satz von Heine-Borel/Fakt|Satz 13.14 (Differentialgeometrie (Osnabrück 2023))]]
ist die $S^1$ als abgeschlossene und beschränkte Teilmenge des $\R^2$
\definitionsverweis {kompakt}{}{,}
das offene Intervall ist dagegen nicht kompakt.

b) Es ist

\mavergleichskettedisp
{\vergleichskette
{ \omega
}
{ =} { hdt
}
{ } {
}
{ } {
}
{ } {
}
}
{}{}{}
mit einer stetigen Funktion
\maabbdisp {h} {]0,1[} {\R
} {.}
Nach
dem [[Hauptsatz der Infinitesimalrechnung/Riemann/Fakt|Hauptsatz der Infinitesimalrechnung]]
besitzt $h$ eine Stammfunktion $H$. In der Sprache der Differentialformen bedeutet dies

\mavergleichskettedisp
{\vergleichskette
{H'
}
{ =} { hdt
}
{ =} { \omega
}
{ } {
}
{ } {
}
}
{}{}{,}
also ist $\omega$ exakt.

c) Es sei
\mathl{S^1 \subset \R^2}{} gegeben, wobei die Koordinaten des $\R^2$ mit
\mathkor {} {u} {und} {v} {}
bezeichnet seien. Wir betrachten die Differentialform

\mavergleichskettedisp
{\vergleichskette
{ \omega
}
{ =} { udv-vdu
}
{ } {
}
{ } {
}
{ } {
}
}
{}{}{}
auf $S^1$. Da $S^1$ eindimensional ist, ist sie geschlossen. Für den Weg
\maabbeledisp {\gamma} {[0, 2 \pi]} {S^1
} { t } { \begin{pmatrix}  \cos  t  \\ \sin  t     \end{pmatrix}
} {,}
ist

\mavergleichskettedisp
{\vergleichskette
{ \int_\gamma \omega
}
{ =} { \int_0^{2 \pi} \gamma^* \omega  dt
}
{ =} { \int_0^{2 \pi}  \cos  t  \cdot \cos  t  +  \sin  t   \sin  t   dt
}
{ =} { \int_0^{2 \pi} 1 dt
}
{ =} { 2 \pi
}
}
{}{}{.}
Nach
[[Mannigfaltigkeiten ohne Rand/Satz von Stokes/Fakt|Korollar 23.3 (Differentialgeometrie (Osnabrück 2023))]]
kann $\omega$ nicht exakt sein.

 }

__NOEDITSECTION__

\inputaufgabeklausurloesung
{4}

{

Es sei

\mavergleichskette
{\vergleichskette
{ U
}
{ \subseteq }{ \R^n
}
{  }{
}
{    }{
}
{   }{
}
}
{}{}{}
\definitionsverweis {offen}{}{}
und
\maabb {f} { U } { \R
} {}
eine zweifach
\definitionsverweis {stetig differenzierbare Funktion}{}{.}
Zeige

\mavergleichskettedisp
{\vergleichskette
{ d(df)
}
{ =} { 0
}
{ } {
}
{ } {
}
{ } {
}
}
{}{}{,}
wobei $d$ die
\definitionsverweis {äußere Ableitung}{}{}
bezeichnet.

}
{

Für eine $1$-Form

\mavergleichskette
{\vergleichskette
{ \omega
}
{ = }{ \sum_{j = 1}^n g_jdx_j
}
{  }{
}
{    }{
}
{   }{
}
}
{}{}{}
ist unter Verwendung von

\mavergleichskette
{\vergleichskette
{ dx_i \wedge dx_j
}
{ = }{ -dx_j \wedge dx_i
}
{  }{
}
{    }{
}
{   }{
}
}
{}{}{}

\mavergleichskettealign
{\vergleichskettealign
{d \omega
}
{ =} { \sum_{j = 1 }^n  d(g_j dx_j)
}
{ =} { \sum_{j = 1 }^n   \left(  \sum_{i = 1}^n { \frac{   \partial  g_j   }{  \partial   x_i   } } dx_i \wedge dx_j  \right)
}
{ =} { \sum _{1 \leq i <j\leq n} \left(  { \frac{   \partial  g_j   }{  \partial   x_i   } } - { \frac{   \partial  g_i   }{  \partial   x_j   } }  \right) dx_i \wedge dx_j
}
{ } {
}
}
{}
{}{.}
Für eine zweimal stetig differenzierbare Funktion $f$ ist

\mavergleichskette
{\vergleichskette
{ df
}
{ = }{ \sum_{j = 1}^n g_j dx_j
}
{  }{
}
{    }{
}
{   }{
}
}
{}{}{}
mit den
\definitionsverweis {partiellen Ableitungen}{}{}

\mavergleichskette
{\vergleichskette
{ g_j
}
{ = }{ { \frac{   \partial   f   }{  \partial   x_j   } }
}
{  }{
}
{    }{
}
{   }{
}
}
{}{}{,}
und daher ist

\mavergleichskette
{\vergleichskette
{ d(df)
}
{ = }{ 0
}
{  }{
}
{    }{
}
{   }{
}
}
{}{}{}
nach
dem [[Differenzierbarkeit/Satz von Schwarz/Fakt|Satz von Schwarz]].

 }

__NOEDITSECTION__

\inputaufgabeklausurloesung
{3}

{

Beschreibe die
\definitionsverweis {Abbildung}{}{}
\maabbeledisp {\psi} { \Complex \setminus \{1\} } { \Complex
} {w} { { \frac{   (w+1)  { \mathrm i}  }{ 1-w  } }
} {}
in reellen Koordinaten.

}
{

Wir schreiben

\mavergleichskette
{\vergleichskette
{ w
}
{ = }{ a+b { \mathrm i}
}
{  }{
}
{    }{
}
{   }{
}
}
{}{}{.}
Dann ist

\mavergleichskettealign
{\vergleichskettealign
{  \psi ( a+b { \mathrm i}   )
}
{ =} { { \frac{   \left(  a+b { \mathrm i} +1  \right) { \mathrm i}      }{  1-  a-b { \mathrm i}      } }
}
{ =} { { \frac{     a { \mathrm i} +{ \mathrm i}  -b  }{  1-  a-b { \mathrm i}      } }
}
{ =} { { \frac{     \left(  a { \mathrm i} +{ \mathrm i}  -b  \right) \left(  1-  a+b { \mathrm i}  \right)  }{  \left(  1-  a-b { \mathrm i}  \right) \left(  1-  a+b { \mathrm i}  \right)      } }
}
{ =} { { \frac{     \left(  a+1 \right) \left(  1-a  \right) { \mathrm i}  -b^2 { \mathrm i} -b \left(  1-  a \right) -b\left(  a+1  \right)  }{  \left(  1-  a \right)^2 +b^2  } }
}
}
{
\vergleichskettefortsetzungalign
{ =} { { \frac{     \left(  1-a ^2-b^2  \right) { \mathrm i}   -2 b  }{  \left(  1-  a \right)^2 +b^2  } }
}
{ } {
}
{ } {}
{ } {}
}
{}{.}
In reellen Koordinaten ist daher

\mavergleichskettedisp
{\vergleichskette
{ \psi \begin{pmatrix}  a  \\b   \end{pmatrix}
}
{ =} { { \frac{   1  }{  -2a+1+a^2+b^2 } } \begin{pmatrix}  -2b \\ 1-a^2-b^2   \end{pmatrix}
}
{ } {
}
{ } {
}
{ } {
}
}
{}{}{.}

 }

__NOEDITSECTION__

\inputaufgabeklausurloesung
{3}

{

Es sei

\mavergleichskette
{\vergleichskette
{ H
}
{ \subseteq }{ \R^n
}
{  }{
}
{    }{
}
{   }{
}
}
{}{}{}
ein
\definitionsverweis {euklidischer Halbraum}{}{}
und

\mavergleichskette
{\vergleichskette
{ P
}
{ \in }{ H
}
{  }{
}
{    }{
}
{   }{
}
}
{}{}{.}
Es gebe eine in $\R^n$ offene Menge $V$ mit

\mavergleichskette
{\vergleichskette
{ P
}
{ \in }{ V
}
{ \subseteq }{ H
}
{    }{
}
{   }{
}
}
{}{}{.}
Zeige, dass
\mathl{P}{} kein Randpunkt von $H$ ist.

}
{

Es sei

\mavergleichskettedisp
{\vergleichskette
{ P
}
{ =} { \begin{pmatrix}  x_1  \\ x_2 \\  \vdots\\ x_n   \end{pmatrix}
}
{ } {
}
{ } {
}
{ } {
}
}
{}{}{.}
Wir können die im $\R^n$ offene Umgebung $V$ durch einen
\zusatzklammer {im $\R^n$} {} {}
offenen Ball
\mathl{U \left(   P , \epsilon  \right)}{} mit

\mavergleichskette
{\vergleichskette
{ P
}
{ \in }{   U \left(   P , \epsilon  \right)
}
{ \subseteq }{  V
}
{ \subseteq   }{ H
}
{   }{
}
}
{}{}{}
ersetzen. Wenn $P$ ein Randpunkt wäre, so wäre

\mavergleichskette
{\vergleichskette
{ x_1
}
{ = }{ 0
}
{  }{
}
{    }{
}
{   }{
}
}
{}{}{.}
Doch dann wäre

\mavergleichskette
{\vergleichskette
{ \begin{pmatrix}  - { \frac{   \epsilon }{  2 } }  \\ x_2 \\  \vdots\\ x_n   \end{pmatrix}
}
{ \in }{ U \left(   P , \epsilon  \right)
}
{  }{
}
{    }{
}
{   }{
}
}
{}{}{,}
aber dieser Punkt gehört nicht zu $H$.

 }

__NOEDITSECTION__

\inputaufgabeklausurloesung
{7}

{

Es sei $M$ eine $n$-dimensionale
\definitionsverweis {differenzierbare Mannigfaltigkeit}{}{}
und

\mavergleichskette
{\vergleichskette
{ P
}
{ \in }{ M
}
{  }{
}
{    }{
}
{   }{
}
}
{}{}{}
ein Punkt. Zeige, dass es eine
\definitionsverweis {Mannigfaltigkeit mit Rand}{}{}
$N$ gibt, deren Rand $\partial N$  diffeomorph zur $(n-1)$-dimensionalen Sphäre $S^{n-1}$ ist und derart, dass
\mathkor {} {M \setminus \{ P \}} {und} {N \setminus \partial N} {}
zueinander diffeomorph sind.

}
{

Es sei

\mavergleichskette
{\vergleichskette
{ P
}
{ \in }{  U
}
{ \subseteq }{  M
}
{    }{
}
{   }{
}
}
{}{}{}
ein offenes Kartengebiet mit der Karte
\maabbdisp {\alpha} { U } {V \subseteq \R^n
} {.}
Wir können davon ausgehen, dass $\alpha(P)$ der Nullpunkt ist und das $V$ der offene Ball mit Radius $3$ ist. Es sei
\maabbdisp {\varphi} { U \left(   0 , 3  \right) \setminus \{0\}  } { U \left(   0 , 3  \right) \setminus B  \left(  0 , 1 \right)
} {}
ein Diffemorphismus, der auf
\mathl{U \left(   0 , 3  \right) \setminus U \left(   0 , 2  \right)}{} die Identität ist und der
\mathl{U \left(   0 , 2  \right) \setminus \{0\}}{} auf
\mathl{U \left(   0 , 2  \right) \setminus B  \left(  0 , 1 \right)}{} abbildet. Eine solche Abbildung erhält man, wenn man mit der Funktion
\maabbdisp {h} {[0,3]} { [0,3]
} {}
mit

\mavergleichskettedisp
{\vergleichskette
{ h(r)
}
{ =} { \begin{cases} 1+ { \frac{   1  }{  4 } } r^2 \text{ für } 0 \leq r \leq 2 \, , \\ r \text{ für } r >  2  \, , \end{cases}
}
{ } {
}
{ } {
}
{ } {
}
}
{}{}{}
die Punkte streckt. Dabei kann man das Bild des Diffeomorphismus $\varphi$, also
\mathl{U \left(   0 , 3  \right) \setminus B  \left(  0 , 1 \right)}{,} als

\mavergleichskettedisp
{\vergleichskette
{  U \left(   0 , 3  \right) \setminus B  \left(  0 , 1 \right)
}
{ \subset} {  U \left(   0 , 3  \right) \setminus U \left(   0 , 1  \right)
}
{ } {
}
{ } {
}
{ } {
}
}
{}{}{}
auffassen, wobei die größere Menge eine Mannigfaltigkeit mit dem Rand $S \left(   0 ,  1  \right)$ ist. Es sei $N$ die Mannigfaltigkeit, die entsteht, wenn man $U$ durch
\mathl{U \left(   0 , 3  \right) \setminus U \left(   0 , 1  \right)}{} ersetzt. Die Sphäre wird dabei zum Rand von $N$. Den Diffeomorphismus $\varphi$ man man zu einem Diffeomorphismus
\maabbdisp {} {M \setminus \{P\} } { N \setminus \partial N
} {}
fortsetzen, da auf dem offenen Übergang
\mathl{U \left(   0 , 3  \right) \setminus U \left(   0 , 2  \right)}{} die Identität vorliegt.

 }

__NOEDITSECTION__

\inputaufgabeklausurloesung
{2}

{

Man gebe eine
\definitionsverweis {kompakte Ausschöpfung}{}{}
für den $\R^d$ an.

}
{

Wir betrachten die
\definitionsverweis {abgeschlossenen Bälle}{}{}

\mathbed {B  \left(  0 , n  \right)} {}
 {n \in \N_+} {}
 {} {} {} {,}
die
\definitionsverweis {kompakt}{}{}
sind und die den $\R^d$ überdecken. Das offene Innere von
\mathl{B  \left(  0 , n+1  \right)}{}
ist der
\definitionsverweis {offene Ball}{}{}

\mathl{U \left(   0,n+1  \right)}{} und wegen

\mavergleichskettedisp
{\vergleichskette
{ B  \left(  0,n \right)
}
{ \subseteq} { U \left(   0,n+1  \right)
}
{ } {
}
{ } {
}
{ } {
}
}
{}{}{}
liegt eine
\definitionsverweis {kompakte Ausschöpfung}{}{}
vor.

 }

__NOEDITSECTION__

\inputaufgabeklausurloesung
{9}

{

Beweise den Brouwerschen Fixpunktsatz.

}
{

Zur Notationsvereinfachung sei

\mavergleichskette
{\vergleichskette
{ r
}
{ = }{ 1
}
{  }{
}
{    }{
}
{   }{
}
}
{}{}{.}  Nehmen wir an, dass es eine fixpunktfreie stetig differenzierbare Abbildung $\psi$ geben würde. Dann ist stets

\mavergleichskettedisp
{\vergleichskette
{ x
}
{ \neq} { \psi(x)
}
{ } {
}
{ } {
}
{ } {
}
}
{}{}{,}
sodass die beiden Punkte eine Gerade definieren. Die Idee ist, mittels dieser Geraden einen
\zusatzklammer {der beiden} {} {}
Durchstoßungspunkt mit der Sphäre als Bildpunkt einer Retraktion auf den Rand zu nehmen. Mit der Hilfsfunktion

\mavergleichskettedisp
{\vergleichskette
{ h(x)
}
{ =} { { \frac{   \psi(x)-x }{  \Vert { \psi (x)-x} \Vert  } }
}
{ } {
}
{ } {
}
{ } {
}
}
{}{}{}
definieren wir eine Abbildung
\maabbdisp {\varphi} { B  \left(  0 , 1 \right) } { \R^n
} {}
durch

\mavergleichskettedisphandlinks
{\vergleichskettedisphandlinks
{ \varphi(x)
}
{ =} { x + \left(  - \left\langle  x  , h(x)  \right\rangle + \sqrt{1 + \left\langle  x  , h(x)  \right\rangle^2 - \Vert { x} \Vert^2  }  \right) \cdot h(x)
}
{ } {
}
{ } {
}
{ } {
}
}
{}{}{.}
Dabei ist der Ausdruck unter der Wurzel positiv. Dies ist bei

\mavergleichskette
{\vergleichskette
{ \Vert { x } \Vert
}
{ < }{ 1
}
{  }{
}
{    }{
}
{   }{
}
}
{}{}{}
klar und bei

\mavergleichskette
{\vergleichskette
{ \Vert { x } \Vert
}
{ = }{ 1
}
{  }{
}
{    }{
}
{   }{
}
}
{}{}{}
liegt ein Punkt auf der Sphäre vor, dessen Verbindungsgerade mit dem Kugelpunkt
\mathl{\psi(x)}{} nicht senkrecht zu $x$ ist
\zusatzklammer {der affine Tangentialraum zu einem Punkt der Sphäre trifft eine Kugel nur in einem Punkt} {} {,}
sodass

\mavergleichskettedisp
{\vergleichskette
{ \left\langle  x  , h(x)  \right\rangle
}
{ \neq} { 0
}
{ } {
}
{ } {
}
{ } {
}
}
{}{}{}
ist. Da die Quadratwurzel und der Betrag außerhalb des Nullpunktes
\definitionsverweis {stetig differenzierbar}{}{}
sind, handelt es sich bei
\mathkor {} {h(x)} {und bei} {\varphi(x)} {}
um stetig differenzierbare Abbildungen. Die Abbildung $\varphi$ bildet nach
[[Brouwersche Fixpunktsatz/Explizite Retraktion/Aufgabe|Aufgabe 23.23 (Differentialgeometrie (Osnabrück 2023))]]
die Kugel auf die Sphäre ab und ihre Einschränkung auf die Sphäre ist die Identität. Damit liegt eine stetig differenzierbare
\definitionsverweis {Retraktion}{}{}
der abgeschlossenen Vollkugel auf ihren Rand vor, was nach
[[Mannigfaltigkeit mit Rand/Stokes/Nichtexistenz von Retraktionen/Fakt|Satz 23.9 (Differentialgeometrie (Osnabrück 2023))]]
nicht sein kann.

 }

__NOEDITSECTION__

\inputaufgabeklausurloesung
{7 (1+2+4)}

{

Wir betrachten auf dem
\definitionsverweis {trivialen Vektorbündel}{}{}

\mathl{\R \times \R^2}{} über $\R$ den
\definitionsverweis {trivialen Zusammenhang}{}{}
$\nabla$.
\aufzaehlungdrei{Zeige, dass

\mavergleichskette
{\vergleichskette
{ f_1(t)
}
{ = }{ \left(  1+ t^2  , \,  t^3                   \right)
}
{  }{
}
{    }{
}
{   }{
}
}
{}{}{}
und

\mavergleichskette
{\vergleichskette
{ f_2(t)
}
{ = }{ \left(  t  , \,   1+t^2                   \right)
}
{  }{
}
{    }{
}
{   }{
}
}
{}{}{}
\definitionsverweis {Basisschnitte}{}{}
des Vektorbündels sind.
}{Drücke die Standardbasisschnitte $e_1,e_2$ als Linearkombination der Basisschnitte $f_1,f_2$ aus.
}{Bestimme die
\definitionsverweis {Christoffelsymbole}{}{}
von $\nabla$ bezüglich dieser Basisschnitte.
}

}
{

\aufzaehlungdrei{Es ist

\mavergleichskettedisp
{\vergleichskette
{ \det  \begin{pmatrix}  1+t^2   &  t^3  \\  t  &  1+t^2  \end{pmatrix}
}
{ =} { 1+2t^2+t^4 -t^4
}
{ =} { 1 +2t^2
}
{ >} {  0
}
{ } {
}
}
{}{}{}
stets positiv, daher sind die beiden Vektoren
\mathkor {} {f_1(t)} {und} {f_2(t)} {}
zu jedem

\mavergleichskette
{\vergleichskette
{ t
}
{ \in }{  \R
}
{  }{
}
{    }{
}
{   }{
}
}
{}{}{}
eine Basis.
}{Aus der Beziehung

\mavergleichskettedisp
{\vergleichskette
{ \begin{pmatrix} f_1  \\f_2   \end{pmatrix}
}
{ =} { \begin{pmatrix}  1+t^2  &   t^3  \\  t  &  1+t^2 \end{pmatrix}  \begin{pmatrix} e_1  \\e_2   \end{pmatrix}
}
{ } {
}
{ } {
}
{ } {
}
}
{}{}{}
folgt

\mavergleichskettedisp
{\vergleichskette
{ \begin{pmatrix} e_1  \\e_2   \end{pmatrix}
}
{ =} { { \frac{   1  }{  1+2t^2 } } \begin{pmatrix}  1+t^2  &   -t^3  \\  -t &  1+t^2 \end{pmatrix}  \begin{pmatrix} f_1  \\f_2   \end{pmatrix}
}
{ } {
}
{ } {
}
{ } {
}
}
{}{}{,}
also

\mavergleichskettedisp
{\vergleichskette
{ e_1
}
{ =} { { \frac{   1+t^2 }{  1+2t^2 } }   f_1 - { \frac{  t^3 }{  1+2t^2 } } f_2
}
{ } {
}
{ } {
}
{ } {
}
}
{}{}{}
und

\mavergleichskettedisp
{\vergleichskette
{ e_2
}
{ =} { -{ \frac{   t  }{  1+2t^2 } }   f_1 + { \frac{   1+t^2 }{  1+2t^2 } } f_2
}
{ } {
}
{ } {
}
{ } {
}
}
{}{}{.}
}{Für den trivialen Zusammenhang ist

\mavergleichskettedisp
{\vergleichskette
{ \nabla_\partial (f)
}
{ =} { \partial f
}
{ } {
}
{ } {
}
{ } {
}
}
{}{}{.}
Somit ist

\mavergleichskettealign
{\vergleichskettealign
{ \nabla_\partial (f_1)
}
{ =} { \nabla_\partial \left(  \left(  1+t^2  \right) e_1 + t^3e_2  \right)
}
{ =} { 2t e_1 + 3t^2 e_2
}
{ =} { 2t  \left(  { \frac{   1+t^2 }{  1+2t^2 } }   f_1 - { \frac{  t^3 }{  1+2t^2 } } f_2   \right)  + 3t^2  \left(  -{ \frac{   t  }{  1+2t^2 } }   f_1 + { \frac{   1+t^2 }{  1+2t^2 } } f_2   \right)
}
{ =} {  \left(   { \frac{   2t+2t^3 }{  1+2t^2 } } - { \frac{   3t^3 }{  1+2t^2 } }  \right) f_1  +  \left(  - { \frac{  t^3 }{  1+2t^2 } }   + { \frac{   3t^2+3t^4 }{  1+2t^2 } }     \right)   f_2
}
}
{
\vergleichskettefortsetzungalign
{ =} { { \frac{   2t - t^3 }{  1+2t^2 } }  f_1  +  { \frac{   3t^2-t^3+3t^4 }{  1+2t^2 } }      f_2
}
{ } {}
{ } {}
{ } {}
}
{}{}
und

\mavergleichskettealign
{\vergleichskettealign
{ \nabla_\partial (f_2)
}
{ =} { \nabla_\partial \left(  t e_1 + \left(  1+t^2  \right) e_2  \right)
}
{ =} { e_1 + 2t  e_2
}
{ =} { { \frac{   1+t^2 }{  1+2t^2 } } f_1 - { \frac{  t^3 }{  1+2t^2 } } f_2 + 2t \left(  -{ \frac{   t  }{  1+2t^2 } } f_1 + { \frac{   1+t^2 }{  1+2t^2 } } f_2   \right)
}
{ =} { \left( { \frac{   1+t^2 }{  1+2t^2 } }  - { \frac{   2t^2 }{  1+2t^2 } }  \right) f_1  + \left(  - { \frac{   t^3  }{  1+2t^2  } } +  { \frac{   2t+2t^3  }{  1+2t^2 } }  \right) f_2
}
}
{
\vergleichskettefortsetzungalign
{ =} { { \frac{   1 - t^2  }{  1+2t^2 } } f_1 + { \frac{   2t   +t^3 }{  1+2t^2 } }      f_2
}
{ } {}
{ } {}
{ } {}
}
{}{.}
Die Koeffizientenfunktionen sind die Christoffelsymbole.
}

 }

 [[Kategorie:Latexseite]]
